\subsection{SUPPLEMENTARY SUBMODULES}

DEFINITION 8. In an A-module E, two submodules \(M_{1}\) , \(M_{2}\) are said to be supplementary if E is the direct sum of \(M_{1}\) and \(M_{2}\) .

Proposition 11 of no. 8 shows that, for \(\mathbf{M}_1\) and \(\mathbf{M}_2\) to be supplementary, it is necessary and sufficient that \(\mathbf{M}_1 + \mathbf{M}_2 = \mathbf{E}\) and \(\mathbf{M}_1 \cap \mathbf{M}_2 = \{0\}\) (cf.~I, §4, no. 9, Proposition 15).

PROPOSITION 13. Let \(\mathbf{M}_1, \mathbf{M}_2\) be two supplementary submodules in an A-module E. The restriction to \(\mathbf{M}_1\) of the canonical mapping \(\mathbf{E} \to \mathbf{E} / \mathbf{M}_2\) is an isomorphism of \(\mathbf{M}_1\) onto \(\mathbf{E} / \mathbf{M}_2\) .

This linear mapping is surjective since \(M_{1} + M_{2} = E\) and it is injective since its kernel is the intersection of \(M_{1}\) and the kernel \(M_{2}\) of \(E \rightarrow E/M_{2}\) and hence reduces to \(\{0\}\) .

COROLLARY. If \(M_{2}\) and \(M_{2}^{\prime}\) are two supplements of the same submodule \(M_{1}\) of E, the set of ordered pairs \((x, x^{\prime}) \in \mathbf{M}_{2} \times \mathbf{M}_{2}^{\prime}\) such that \(x - x^{\prime} \in M_{1}\) is the graph of an isomorphism of \(M_{2}\) onto \(M_{2}^{\prime}\) .

It is immediate that it is the graph of the composite isomorphism \(M_{2} \rightarrow E/M_{1} \rightarrow M_{2}'\) .

DEFINITION 9. A submodule M of an A-module E is called a direct factor of E if it has a supplementary submodule in E.

When this is so, E/M is isomorphic to a supplement of M (Proposition 13).

A submodule does not necessarily admit a supplement (Exercise 11). When a submodule is a direct factor, it has in general several distinct supplements; these supplements are however canonically isomorphic to one another (Corollary to Proposition 13).

PROPOSITION 14. For a submodule M of a module E to be a direct factor, it is necessary and sufficient that there exist a projector of E whose image is M or a projector of E whose kernel is M.

This follows immediately from no. 8, Proposition 12 and Corollary.

PROPOSITION 15. Given an exact sequence of A-modules

\[
0 \longrightarrow \mathrm{E} \stackrel {f} {\longrightarrow} \mathrm{F} \stackrel {g} {\longrightarrow} \mathrm{G} \longrightarrow 0\tag{33}
\]

the following propositions are equivalent:

\begin{enumerate}
\def\labelenumi{(\alph{enumi})}
\item
  The submodule \(f(\mathbf{E})\) of \(\mathbf{F}\) is a direct factor.
\item
  There exists a linear retraction \(r: \mathbf{F} \to \mathbf{E}\) associated with \(f\) (Set Theory, II, § 3, no. 8, Definition 11).
\item
  There exists a linear section \(s: \mathbf{G} \to \mathbf{F}\) associated with \(g\) (Set Theory, II, § 3, no. 8, Definition 11).
\end{enumerate}

When this is so, \(f + s: \mathbf{E} \oplus \mathbf{G} \to \mathbf{F}\) is an isomorphism.

If there exists a projector \(e\) in \(\operatorname{End}(\mathbf{F})\) such that \(e(\mathbf{F}) = f(\mathbf{E})\) , the homomorphism \(f^{-1} \circ e: \mathbf{F} \to \mathbf{E}\) is a linear retraction associated with \(f\) ; conversely, if there exists such a retraction \(r\) , it is immediate that \(f \circ r\) is a projector in \(\mathbf{F}\) whose image is \(f(\mathbf{E})\) , hence (a) and (b) are equivalent (Proposition 14). If \(f(\mathbf{E})\) admits a supplement \(\mathbf{E}'\) in \(\mathbf{F}\) and \(j: \mathbf{E}' \to \mathbf{F}\) is the canonical injection, \(g \circ j\) is an isomorphism of \(\mathbf{E}'\) onto \(\mathbf{G}\) and the inverse isomorphism, considered as a mapping of \(\mathbf{G}\) into \(\mathbf{F}\) , is a linear section associated with \(g\) . Conversely, if such a section \(s\) exists, \(s \circ g\) is a projector in \(\mathbf{F}\) whose kernel is \(f(\mathbf{E})\) , hence (a) and (c) are equivalent (Proposition 14). Moreover \(s\) is a bijection of \(\mathbf{G}\) onto \(s(\mathbf{G})\) and as \(s(\mathbf{G})\) is supplementary to \(f(\mathbf{E}), f + s\) is an isomorphism.

Note that being given r (resp. s) is equivalent to being given a supplement of \(f(\mathrm{E})\) in F, namely the kernel of r (resp. image of s).

When the exact sequence (33) satisfies the conditions of Proposition 15, it is said to split or \((\mathbf{F}, f, g)\) is said to be a trivial extension of \(\mathbf{G}\) by \(\mathbf{E}\) (I, § 6, no. 1).

COROLLARY 1. Let \(u: \mathbf{E} \to \mathbf{F}\) be a linear mapping. For there to exist a linear mapping \(v: \mathbf{F} \to \mathbf{E}\) such that \(u \circ v = 1_{\mathbf{F}}\) (the case where \(u\) is said to be right invertible and \(v\) is said to be the right inverse of \(u\) ), it is necessary and sufficient that \(u\) be surjective and that its kernel be a direct factor in \(\mathbf{E}\) . The submodule \(\operatorname{Im}(v)\) of \(\mathbf{E}\) is then a supplement of \(\operatorname{Ker}(u)\) .

It is obviously necessary that \(u\) be surjective; as \(v\) is then a section associated with \(u\) , the conclusion follows from Proposition 15.

COROLLARY 2. Let \(u: \mathbf{E} \to \mathbf{F}\) be a linear mapping. For there to exist a linear mapping \(v: \mathbf{F} \to \mathbf{E}\) such that \(v \circ u = 1_{\mathbf{E}}\) (the case where \(u\) is said to be left invertible and \(v\) is said to be the left inverse of \(u\) ), it is necessary and sufficient that \(u\) be injective and that its image be a direct factor in \(\mathbf{F}\) . The submodule \(\operatorname{Ker}(v)\) of \(\mathbf{F}\) is then a supplement of \(\operatorname{Im}(u)\) .

It is obviously necessary that u be injective; as v is then a retraction associated with u, the conclusion also follows from Proposition 15.

Remarks. (1) Let M, N be two supplementary submodules in an A-module E, p, q the projectors of E onto M and N respectively, corresponding to the decomposition of E as a direct sum of M and N. It is known (no. 6, Corollary 1 to Proposition 6) that, for every A-module F, the mapping \((u, v) \mapsto u \circ p + v \circ q\) is an isomorphism of

\[
\operatorname{Hom} _ {\mathbf {A}} (\mathbf {M}, \mathbf {F}) \oplus \operatorname{Hom} _ {\mathbf {A}} (\mathbf {N}, \mathbf {F})
\]

onto \(\operatorname{Hom}_{\mathbf{A}}(\mathbf{E},\mathbf{F})\) . The image of \(\operatorname{Hom}_{\mathbf{A}}(\mathbf{M},\mathbf{F})\) under this isomorphism is the set of linear mappings \(w:\mathrm{E}\to \mathrm{F}\) such that \(w(x) = 0\) for all \(x\in N\) .

\begin{enumerate}
\def\labelenumi{(\arabic{enumi})}
\setcounter{enumi}{1}
\tightlist
\item
  If M, N are two submodules of E such that \(M \cap N\) is a direct factor of M and N, then \(M \cap N\) is also a direct factor of \(M + N\) : if P (resp. Q) is a supplement of \(M \cap N\) in M (resp. N), \(M + N\) is the direct sum of \(M \cap N\) , P and Q, as is immediately verified.
\end{enumerate}

